\chapter{Fundamentos da Definição de Dados}
\label{cap:fundamentos_definicao_dados}

A linguagem de definição de dados, conhecida como DDL (\textit{Data Definition Language}), reúne os comandos responsáveis por criar, organizar, ajustar e remover estruturas de um banco de dados. Enquanto a Aula 10 apresentou a SQL como linguagem geral de implementação, esta aula concentra-se no conjunto de comandos que materializam o projeto lógico em estruturas concretas.

\section{O papel da DDL no desenvolvimento do banco}

No processo de construção de um banco de dados, a DDL cumpre papel estruturante. É ela que traduz entidades, atributos, chaves e relacionamentos para formas efetivamente compreendidas pelo SGBD. Sem DDL, o modelo permaneceria apenas no plano do projeto.

Isso significa que a DDL é responsável por:

\begin{itemize}
    \item criar bancos e schemas;
    \item criar tabelas;
    \item definir colunas e tipos;
    \item estabelecer restrições;
    \item remover ou alterar estruturas.
\end{itemize}

\section{Comandos centrais da DDL}

Entre os comandos mais importantes da DDL, destacam-se:

\begin{itemize}
    \item \texttt{CREATE DATABASE}
    \item \texttt{CREATE SCHEMA}
    \item \texttt{CREATE TABLE}
    \item \texttt{ALTER TABLE}
    \item \texttt{DROP TABLE}
\end{itemize}

Nesta aula, o foco principal recairá sobre a criação inicial de estruturas. Alterações posteriores serão aprofundadas na Aula 12.

\section{Criação de banco de dados}

\begin{sqlcode}
CREATE DATABASE escola;
\end{sqlcode}

A criação de um banco representa o ponto de partida da implementação. Em um ambiente acadêmico, isso pode corresponder à base que armazenará as informações de alunos, cursos, professores, turmas e matrículas.

\section{Criação de schema}

\begin{sqlcode}
CREATE SCHEMA academico;
\end{sqlcode}

O schema organiza objetos logicamente. Em projetos maiores, esse recurso ajuda a separar módulos e responsabilidades.

Exemplo:
\begin{itemize}
    \item schema \texttt{academico};
    \item schema \texttt{financeiro};
    \item schema \texttt{biblioteca}.
\end{itemize}

\section{Criação de tabelas}

A criação de tabelas é o núcleo da DDL. É nesse momento que se definem as estruturas que receberão os dados da aplicação.

\subsection{Estrutura mínima}

\begin{sqlcode}
CREATE TABLE academico.aluno (
    aluno_id SERIAL PRIMARY KEY,
    nome VARCHAR(100) NOT NULL,
    idade INTEGER,
    curso VARCHAR(100) NOT NULL
);
\end{sqlcode}

Nesse exemplo, aparecem decisões fundamentais:
\begin{itemize}
    \item uma coluna identificadora;
    \item colunas textuais;
    \item uma coluna numérica;
    \item obrigatoriedade em campos essenciais.
\end{itemize}

\section{Tipos de dados}

A escolha do tipo de dado é uma decisão de projeto físico. Ela determina:
\begin{itemize}
    \item a natureza dos valores aceitos;
    \item o espaço de armazenamento;
    \item a semântica da coluna;
    \item algumas possibilidades de validação.
\end{itemize}

\subsection{Tipos numéricos}

\begin{itemize}
    \item \texttt{INTEGER}: inteiros;
    \item \texttt{BIGINT}: inteiros maiores;
    \item \texttt{NUMERIC(p,s)}: valores exatos com precisão e escala;
    \item \texttt{REAL} e \texttt{DOUBLE PRECISION}: valores aproximados.
\end{itemize}

\subsection{Tipos textuais}

\begin{itemize}
    \item \texttt{CHAR(n)}: tamanho fixo;
    \item \texttt{VARCHAR(n)}: tamanho variável com limite;
    \item \texttt{TEXT}: texto de tamanho amplo.
\end{itemize}

\subsection{Tipos temporais}

\begin{itemize}
    \item \texttt{DATE}: data;
    \item \texttt{TIME}: hora;
    \item \texttt{TIMESTAMP}: data e hora.
\end{itemize}

\subsection{Tipos úteis no PostgreSQL}

\begin{itemize}
    \item \texttt{SERIAL}: sequência autoincrementável;
    \item \texttt{JSON}: estrutura semiestruturada;
    \item \texttt{ARRAY}: coleção homogênea de valores;
    \item \texttt{daterange}: intervalo de datas.
\end{itemize}

\section{Restrições fundamentais}

As restrições ajudam a preservar a integridade do banco e a aproximar a implementação das regras do domínio.

\subsection{NOT NULL}

Indica que a coluna não pode receber valor nulo.

\begin{sqlcode}
nome VARCHAR(100) NOT NULL
\end{sqlcode}

\subsection{PRIMARY KEY}

Identifica unicamente cada linha da tabela.

\begin{sqlcode}
curso_id SERIAL PRIMARY KEY
\end{sqlcode}

\subsection{FOREIGN KEY}

Estabelece vínculo com outra tabela.

\begin{sqlcode}
FOREIGN KEY (curso_id) REFERENCES curso(curso_id)
\end{sqlcode}

\subsection{UNIQUE}

Impede repetição de valores em determinada coluna ou conjunto de colunas.

\subsection{CHECK}

Permite impor restrições lógicas.

\subsection{DEFAULT}

Define valores padrão.

\section{Chave primária simples e composta}

\subsection{Chave simples}

\begin{sqlcode}
CREATE TABLE curso (
    curso_id SERIAL PRIMARY KEY,
    nome VARCHAR(100) NOT NULL
);
\end{sqlcode}

\subsection{Chave composta}

\begin{sqlcode}
CREATE TABLE inscricao (
    aluno_id INTEGER NOT NULL,
    disciplina_id INTEGER NOT NULL,
    semestre VARCHAR(10) NOT NULL,
    PRIMARY KEY (aluno_id, disciplina_id, semestre)
);
\end{sqlcode}

A chave composta é útil quando a identificação depende da combinação de múltiplos atributos.

\section{Chave estrangeira e implementação de relacionamentos}

A chave estrangeira é o mecanismo relacional que traduz, em tabelas, os vínculos identificados no modelo lógico.

\begin{sqlcode}
CREATE TABLE curso (
    curso_id SERIAL PRIMARY KEY,
    nome VARCHAR(100) NOT NULL
);
\end{sqlcode}

\begin{sqlcode}
CREATE TABLE aluno (
    aluno_id SERIAL PRIMARY KEY,
    nome VARCHAR(100) NOT NULL,
    curso_id INTEGER NOT NULL,
    FOREIGN KEY (curso_id) REFERENCES curso(curso_id)
);
\end{sqlcode}

Nesse caso, um aluno passa a estar associado a um curso já existente.

\section{Exemplo acadêmico mais completo}

\begin{sqlcode}
CREATE TABLE professor (
    professor_id SERIAL PRIMARY KEY,
    nome VARCHAR(120) NOT NULL,
    titulacao VARCHAR(50) NOT NULL
);
\end{sqlcode}

\begin{sqlcode}
CREATE TABLE disciplina (
    disciplina_id SERIAL PRIMARY KEY,
    nome VARCHAR(120) NOT NULL,
    carga_horaria INTEGER NOT NULL
);
\end{sqlcode}

\begin{sqlcode}
CREATE TABLE turma (
    turma_id SERIAL PRIMARY KEY,
    disciplina_id INTEGER NOT NULL,
    professor_id INTEGER NOT NULL,
    semestre VARCHAR(10) NOT NULL,
    FOREIGN KEY (disciplina_id) REFERENCES disciplina(disciplina_id),
    FOREIGN KEY (professor_id) REFERENCES professor(professor_id)
);
\end{sqlcode}

Esse exemplo mostra como a DDL transforma estruturas lógicas em tabelas relacionadas.

\section{Exemplo comercial}

\begin{sqlcode}
CREATE TABLE cliente (
    cliente_id SERIAL PRIMARY KEY,
    nome VARCHAR(100) NOT NULL,
    email VARCHAR(120) UNIQUE
);
\end{sqlcode}

\begin{sqlcode}
CREATE TABLE pedido (
    pedido_id SERIAL PRIMARY KEY,
    cliente_id INTEGER NOT NULL,
    data_pedido DATE NOT NULL,
    valor_total NUMERIC(10,2) NOT NULL,
    FOREIGN KEY (cliente_id) REFERENCES cliente(cliente_id)
);
\end{sqlcode}

\section{Exemplo de recursos humanos}

\begin{sqlcode}
CREATE TABLE departamento (
    departamento_id SERIAL PRIMARY KEY,
    nome VARCHAR(100) NOT NULL
);
\end{sqlcode}

\begin{sqlcode}
CREATE TABLE funcionario (
    funcionario_id SERIAL PRIMARY KEY,
    nome VARCHAR(120) NOT NULL,
    salario NUMERIC(10,2) NOT NULL,
    departamento_id INTEGER NOT NULL,
    FOREIGN KEY (departamento_id) REFERENCES departamento(departamento_id)
);
\end{sqlcode}

\section{Escolhas de projeto físico}

Ao criar estruturas em SQL, o desenvolvedor precisa tomar decisões que não aparecem com esse nível de detalhe no modelo conceitual. Entre elas:

\begin{itemize}
    \item qual tipo de dado adotar;
    \item qual será a chave primária;
    \item que colunas devem ser obrigatórias;
    \item quais colunas devem ter unicidade;
    \item como representar relacionamentos.
\end{itemize}

Essas decisões exigem domínio técnico e compreensão do contexto do sistema.

\section{Boas práticas em DDL}

\begin{itemize}
    \item usar nomes descritivos;
    \item manter padrão de nomenclatura;
    \item definir chaves primárias explicitamente;
    \item evitar tipos exageradamente genéricos quando há alternativa mais precisa;
    \item registrar restrições relevantes desde a criação da tabela.
\end{itemize}

\begin{quadro_dica}
Uma modelagem física bem construída reduz retrabalho e evita que regras essenciais fiquem dependentes apenas da aplicação.
\end{quadro_dica}

\section{Exercícios básicos}

\begin{enumerate}
    \item Crie uma tabela \texttt{professor} com identificador, nome, titulação e e-mail.
    \item Crie uma tabela \texttt{disciplina} com identificador, nome e carga horária.
    \item Crie uma tabela \texttt{turma} relacionando professor e disciplina.
    \item Escreva um exemplo de chave composta.
\end{enumerate}

\section{Exercícios intermediários}

\begin{enumerate}
    \item Compare \texttt{CHAR}, \texttt{VARCHAR} e \texttt{TEXT}.
    \item Explique em quais situações uma chave composta pode ser mais adequada que uma chave simples.
    \item Modele um pequeno cenário comercial com cliente e pedido.
    \item Modele um pequeno cenário de recursos humanos com departamento e funcionário.
\end{enumerate}

\section{Exercício integrador}

Projete, em SQL DDL, a base de um sistema de biblioteca com as tabelas:
\begin{itemize}
    \item \texttt{livro}
    \item \texttt{autor}
    \item \texttt{categoria}
    \item \texttt{leitor}
    \item \texttt{emprestimo}
\end{itemize}

Para cada tabela, defina:
\begin{enumerate}
    \item chave primária;
    \item tipos de dados;
    \item colunas obrigatórias;
    \item relacionamentos por chave estrangeira.
\end{enumerate}

\section{Síntese da aula}

Nesta aula, a DDL foi apresentada como instrumento central da implementação física de bancos de dados. Foram discutidos criação de banco, schema e tabelas, tipos de dados, chaves simples e compostas, chaves estrangeiras e boas práticas de projeto físico. A próxima aula aprofundará alterações estruturais, restrições adicionais e comportamentos de integridade referencial.